\documentclass[11pt]{article}
\usepackage[T1]{fontenc}
\usepackage{graphicx,hyperref}
\usepackage{titling}

\usepackage{xcolor}
\usepackage{amsmath,amssymb}
\usepackage[normalem]{ulem}
\usepackage{fancyvrb}
\usepackage[most]{tcolorbox}

%%%% inspired by classicthesis.sty
\usepackage{newpxtext,newpxmath}
\usepackage{microtype}
\usepackage{titlesec}
\titleformat{\chapter}[display]
  {\relax}{\large\scshape\MakeUppercase{\chaptertitlename}~\thechapter}{0pt}{\large\scshape}
\titleformat{\section}{\large\scshape}{\thesection}{1em}{}
\titleformat{\subsection}{\normalsize\itshape}{\thesubsection}{1em}{}
\titleformat{\subsubsection}{\small\normalfont}{\thesubsubsection}{1em}{}


% Source code box (light gray background, left accent bar)
\newtcolorbox{srccode}{
  enhanced, breakable,
  colback=gray!5, colframe=gray!40,
  left=4pt, right=4pt, top=2pt, bottom=2pt,
  boxrule=0pt, leftrule=3pt,
  arc=0pt, outer arc=0pt
}

% Output box (white background, thin border)
\newtcolorbox{outcode}{
  enhanced, breakable,
  colback=white, colframe=gray!50,
  left=4pt, right=4pt, top=2pt, bottom=2pt,
  boxrule=0.5pt,
  arc=0pt, outer arc=0pt,
  fontupper=\small
}

\title{\emph{Calepin} Tutorial (PDF)}

\author{Vincent Arel-Bundock\textsuperscript{1}\textsuperscript{*} \and Norah Jones\textsuperscript{2,3}}
\newcommand{\affiliationblock}{\begin{center}\small \textsuperscript{1} Department of Political Science, Université de Montréal, Montreal, Canada \\
\textsuperscript{2} Machine Learning, Carnegie Mellon University, Pittsburgh, PA, USA \\
\textsuperscript{3} University of Chicago, Chicago, IL, USA\end{center}}
\date{2026-03-19}
\begin{document}
\maketitle
\ifdefined\affiliationblock\affiliationblock\fi


\begin{abstract}
This document renders the \emph{Calepin} tutorial to PDF via LaTeX. It demonstrates
scholarly front matter (authors, affiliations, ORCID), the appendix, and the
\texttt{include} shortcode for pulling in content from another file.
\end{abstract}



\definecolor{hl1}{RGB}{50,50,50}
\definecolor{hl2}{RGB}{167,29,93}
\definecolor{hl3}{RGB}{0,134,179}
\definecolor{hl4}{RGB}{98,163,92}
\definecolor{hl5}{RGB}{24,54,145}
This page reviews some of the basics of Quarto and Markdown syntax supported by \emph{Calepin}. It covers document structure, code chunks, figures, math, cross-references, citations, and other features that work out of the box.

\section*{Getting started}
\label{getting-started}

\emph{Calepin} renders \texttt{.qmd} documents to HTML, LaTeX, Typst, and Markdown. It executes R code chunks, captures output and plots, resolves citations and cross-references, and wraps everything in a template.

To render this tutorial:



\begin{srccode}
\begin{Verbatim}[commandchars=\\\{\}]
\textcolor{hl1}{calepin tutorial.qmd}
\end{Verbatim}
\end{srccode}

\section*{Markdown basics}
\label{markdown-basics}

\emph{Calepin} supports standard CommonMark and GitHub Flavored Markdown.

Lorem ipsum dolor sit amet, \emph{consectetur} adipiscing elit.
Sed do \textbf{eiusmod tempor} incididunt ut labore.

Ordered list:

\begin{enumerate}
\item First item


\item Second item


\item Third item


\end{enumerate}

Unordered list:

\begin{itemize}
\item Alpha


\item Beta


\item Gamma


\end{itemize}

A \href{https://example.com}{hyperlink} and some \texttt{inline code}.

\begin{quote}
A blockquote: Ut enim ad minim veniam, quis nostrud
exercitation ullamco laboris.

\end{quote}

Text with a footnote\footnotemark[1] and another\footnotemark[2].

Superscript: E = mc\textsuperscript{2}. Strikethrough: \sout{deleted}.

En dash: -- and em dash: ---

\section*{Code chunks}
\label{code-chunks}

Code chunks are fenced with triple backticks and a \texttt{\{r\}} header. Options go inside the chunk as \texttt{\#\textbar{}} pipe comments.

\footnotetext[1]{This is the first footnote.

}\footnotetext[2]{This is the second footnote.

}

\begin{srccode}
\begin{Verbatim}[commandchars=\\\{\}]
\textcolor{hl1}{x }\textcolor{hl2}{<-}\textcolor{hl1}{ }\textcolor{hl3}{1}\textcolor{hl2}{:}\textcolor{hl3}{10}
\textcolor{hl4}{mean}\textcolor{hl1}{(x)}
\end{Verbatim}
\end{srccode}

\begin{outcode}
\begin{verbatim}
[1] 5.5
\end{verbatim}
\end{outcode}

\section*{Inline code}
\label{inline-code}

R expressions can be evaluated inline. The mean of 1:10 is 5.5 and pi is approximately 3.14.

\section*{Figures}
\label{figures}

Code chunks that produce plots automatically generate figures. Use \texttt{fig-cap} for a caption and a label for cross-referencing.



\begin{srccode}
\begin{Verbatim}[commandchars=\\\{\}]
\textcolor{hl4}{plot}\textcolor{hl1}{(}
\textcolor{hl1}{    }\textcolor{hl3}{1}\textcolor{hl2}{:}\textcolor{hl3}{10}\textcolor{hl1}{,}
\textcolor{hl1}{    (}\textcolor{hl3}{1}\textcolor{hl2}{:}\textcolor{hl3}{10}\textcolor{hl1}{)}\textcolor{hl2}{^}\textcolor{hl3}{2}\textcolor{hl1}{,}
\textcolor{hl1}{    main }\textcolor{hl2}{=}\textcolor{hl1}{ }\textcolor{hl5}{"Square Numbers"}\textcolor{hl1}{,}
\textcolor{hl1}{    xlab }\textcolor{hl2}{=}\textcolor{hl1}{ }\textcolor{hl5}{"x"}\textcolor{hl1}{,}
\textcolor{hl1}{    ylab }\textcolor{hl2}{=}\textcolor{hl1}{ }\textcolor{hl5}{"x^2"}
\textcolor{hl1}{)}
\end{Verbatim}
\end{srccode}

\begin{figure}
\centering
\includegraphics[width=0.60\textwidth]{pdf_files/scatter-1.png}
\caption{A simple scatter plot}
\label{fig-scatter}
\end{figure}

\section*{Figure layouts}
\label{figure-layouts}

Figures can be arranged in grids using \texttt{layout-ncol}. The following example places three plots side by side with individual subcaptions and an overall caption.



\begin{figure}
\centering
\begin{minipage}[t]{0.320\textwidth}
\centering
\includegraphics[width=\textwidth]{pdf_files/line-1.png}
\caption{Linear}
\label{fig-line}
\end{minipage}\hfill
\begin{minipage}[t]{0.320\textwidth}
\centering
\includegraphics[width=\textwidth]{pdf_files/quad-1.png}
\caption{Quadratic}
\label{fig-quad}
\end{minipage}\hfill
\begin{minipage}[t]{0.320\textwidth}
\centering
\includegraphics[width=\textwidth]{pdf_files/log-1.png}
\caption{Logarithmic}
\label{fig-log}
\end{minipage}
\caption{Three growth patterns compared.}
\label{fig-panels}
\end{figure}

See \hyperref[fig-panels]{Figure 5} for the result. The \texttt{layout} attribute supports custom grids with relative widths (e.g., \texttt{layout=\textquotedbl{}{[}{[}1,1{]},{[}1{]}{]}\textquotedbl{}}) and negative values for spacing gaps.

\section*{Cross-references}
\label{cross-references}

\emph{Calepin} resolves references to figures, sections, and theorem environments.

\begin{itemize}
\item \hyperref[fig-scatter]{Figure 1} renders as a linked reference


\item Section \ref{sec-figures} links to the Figures section


\item \hyperref[fig-scatter]{1} suppresses the type prefix


\end{itemize}

Sections get automatic ids from their heading text.

\section*{Math}
\label{math}

We can represent dollar amounts using a single dollar sign like 24\$ or enclose math between two dollar signs: $a^2 + b^2 = c^2$.

Display math uses double dollars:

$$
\int_0^1 x^2 \, dx = \frac{1}{3}
$$

\subsection*{Theorem environments}
\label{theorem-environments}

Fenced divs with theorem-type classes get automatic numbering and can be cross-referenced.



\par\medskip\noindent\textbf{Theorem 1.} \label{thm-pythagoras} \itshape{} In a right triangle, the square of the hypotenuse equals the sum of the squares of the other two sides: $a^2 + b^2 = c^2$.

\upshape\medskip

\par\medskip\noindent\textit{Proof.} Let a right triangle have legs $a$, $b$ and hypotenuse $c$. By constructing squares on each side and comparing areas, we obtain $a^2 + b^2 = c^2$.

 \hfill$\square$\medskip

By \hyperref[thm-pythagoras]{Theorem 1}, the relationship is fundamental to Euclidean geometry.

The supported theorem types are: \texttt{theorem}, \texttt{lemma}, \texttt{corollary}, \texttt{proposition}, \texttt{conjecture}, \texttt{definition}, \texttt{example}, \texttt{exercise}, \texttt{solution}, \texttt{remark}, and \texttt{algorithm}. Each type maintains its own counter.

\section*{Callouts}
\label{callouts}

Callout divs highlight important information. Five types are available: \texttt{callout-note}, \texttt{callout-warning}, \texttt{callout-tip}, \texttt{callout-caution}, and \texttt{callout-important}. They can be nested.



\begin{tcolorbox}[enhanced, breakable, colback=blue!5, colframe=blue!40,
  left=4pt, right=4pt, top=4pt, bottom=4pt,
  boxrule=0pt, leftrule=3pt, arc=0pt, outer arc=0pt,
  title=Note, fonttitle=\bfseries]
This is a note callout. Use it for supplementary information.



\begin{tcolorbox}[enhanced, breakable, colback=yellow!10, colframe=yellow!60!black,
  left=4pt, right=4pt, top=4pt, bottom=4pt,
  boxrule=0pt, leftrule=3pt, arc=0pt, outer arc=0pt,
  title=Warning, fonttitle=\bfseries]
Callouts can be nested. This warning appears inside the note above.


\end{tcolorbox}
\end{tcolorbox}

\section*{Conditional content}
\label{conditional-content}

Fenced divs with \texttt{.content-visible} or \texttt{.content-hidden} classes control which content appears in each output format.



Format aliases are supported: \texttt{latex} and \texttt{pdf} both match LaTeX output, \texttt{markdown} and \texttt{md} both match Markdown.

\section*{Raw blocks}
\label{raw-blocks}

Raw blocks inject format-specific markup that passes through verbatim when the output format matches and is dropped otherwise.



\section*{Shortcodes}
\label{shortcodes}

Shortcodes are inline directives written as \texttt{\{\{\textless{} name args \textgreater{}\}\}}. They are processed during rendering and replaced with format-aware output.

\subsection*{pagebreak}
\label{pagebreak}

Insert a native page break. Produces format-appropriate output: \texttt{\textless{}div style=\textquotedbl{}page-break-after: always;\textquotedbl{}\textgreater{}} in HTML, \texttt{\textbackslash{}newpage\{\}} in LaTeX, \texttt{\#pagebreak()} in Typst.

Page one content.



Page two content.

\subsection*{meta}
\label{meta}

Print a value from the document's YAML front matter.

The title of this document is ``\emph{Calepin} Tutorial (PDF)''.

This works with standard fields (\texttt{title}, \texttt{subtitle}, \texttt{author}, \texttt{date}, \texttt{abstract}) and custom metadata.

\subsection*{env}
\label{env}

Print a system environment variable.

The \texttt{SHELL} environment variable on the local machine is: /bin/zsh.

If the variable is not set, the shortcode produces an empty string.

\section*{Citations}
\label{citations}

\emph{Calepin} processes citations from BibTeX files listed in the YAML front matter under \texttt{bibliography}. Use \texttt{@key} for author-year citations and \texttt{{[}-@key{]}} for year-only.

As Arel-Bundock et al. (2026) show, political science research is underpowered. This finding has important implications 2026.

\section*{Tabsets}
\label{tabsets}

Tabsets organize content into switchable tabs. In HTML, readers click to switch tabs. In LaTeX and Typst, tabs render as regular sections.



\subsection*{R}
\label{r}



\begin{srccode}
\begin{Verbatim}[commandchars=\\\{\}]
\textcolor{hl1}{x }\textcolor{hl2}{<-}\textcolor{hl1}{ }\textcolor{hl4}{c}\textcolor{hl1}{(}\textcolor{hl3}{1}\textcolor{hl1}{, }\textcolor{hl3}{2}\textcolor{hl1}{, }\textcolor{hl3}{3}\textcolor{hl1}{)}
\textcolor{hl4}{mean}\textcolor{hl1}{(x)}
\end{Verbatim}
\end{srccode}

\begin{outcode}
\begin{verbatim}
[1] 2
\end{verbatim}
\end{outcode}

\subsection*{Python}
\label{python}



\begin{srccode}
\begin{Verbatim}[commandchars=\\\{\}]
\textcolor{hl1}{x }\textcolor{hl2}{=}\textcolor{hl1}{ [}\textcolor{hl3}{1}\textcolor{hl1}{, }\textcolor{hl3}{2}\textcolor{hl1}{, }\textcolor{hl3}{3}\textcolor{hl1}{]}
\textcolor{hl4}{sum}\textcolor{hl1}{(x) }\textcolor{hl2}{/}\textcolor{hl1}{ }\textcolor{hl4}{len}\textcolor{hl1}{(x)}
\end{Verbatim}
\end{srccode}

\subsection*{Julia}
\label{julia}



\begin{srccode}
\begin{Verbatim}[commandchars=\\\{\}]
\textcolor{hl1}{x = [1, 2, 3]
}
\textcolor{hl1}{sum(x) / length(x)}
\end{Verbatim}
\end{srccode}

Tabsets with the same \texttt{group} attribute switch together:



\subsection*{R}
\label{r}

R version of example 2.

\subsection*{Python}
\label{python}

Python version of example 2.



\subsection*{R}
\label{r}

R content stays synced with the tabset above.

\subsection*{Python}
\label{python}

Python content stays synced too.



\section*{Line blocks}
\label{line-blocks}

Line blocks preserve line breaks and leading spaces. Prefix each line with \texttt{\textbar{}}.

The limerick packs laughs anatomical\\
Into space that is quite economical.\\
~~~But the good ones I've seen\\
~~~So seldom are clean\\
And the clean ones so seldom are comical.

\section*{Footnotes}
\label{footnotes}

There are two ways to create footnotes. Named references use \texttt{{[}\^{}label{]}} at the point of reference and \texttt{{[}\^{}label{]}: text} for the definition elsewhere in the document. Inline footnotes use \texttt{\^{}{[}text{]}} directly without a separate definition.

Named footnote\footnotemark[1] and inline footnote\footnotemark[2].

\section*{Numbered sections}
\label{numbered-sections}

Add \texttt{number-sections: true} to the YAML front matter to automatically number all section headings.

\section*{The \texttt{.hidden} div}
\label{the-hidden-div}

The \texttt{.hidden} div executes its content but produces no output.

\footnotetext[1]{This is defined separately.

}\footnotetext[2]{This is defined right here.

}

The value computed silently: \_\_CALEPIN\_186d\_189e69e21271a060\_\_\_ERROR:object `secret' not found.



\section*{References}
\label{references}

Vincent Arel-Bundock, Ryan C Briggs, Hristos Doucouliagos, Marco M Aviña, Tom D Stanley. 2026. Quantitative political science research is greatly underpowered. The Journal of politics



\appendix
\section*{Appendix}
\subsection*{Reuse}
Creative Commons Attribution 4.0 (\url{https://creativecommons.org/licenses/by/4.0/})
\subsection*{Citation}
Vincent Arel-Bundock, Norah Jones. "*Calepin* Tutorial (PDF)". \emph{Journal of Reproducible Research}. 1(1). (2026-03). DOI: \url{https://doi.org/10.1234/calepin.2026}
\subsection*{Copyright}
Copyright 2026 Vincent Arel-Bundock
\subsection*{Funding}
\begin{itemize}
\item Social Sciences and Humanities Research Council, Grant \#123456, Vincent Arel-Bundock
\end{itemize}


\end{document}
